% ======================================================================
% Chapter 8 -- Case: Credit Card Fraud Review
% Designing Agentic AI Solutions, Chris Kuo
% Part A: brief, KPIs, designs, trade-offs
% ======================================================================

\chapter{Case: Credit Card Fraud Review}
\label{ch:case-fraud}

\begin{caseglance}
\begin{tabularx}{\linewidth}{@{}>{\bfseries}p{3.4cm}X@{}}
Company & Northgate Card Services, a card issuer with 3.2 million credit cards \\
Process & Reviewing card-not-present transactions flagged as possible fraud \\
Starting point & Original design (Part III) \\
Technical focus & Evaluation \\
Business skill & Asymmetric error costs \\
Patterns & Router; approval gate; reflection (for explanations); event-triggered agent \\
Deliverables & Design tables, threshold and capacity model, scoring server, explanation checks, evaluation report, memo
\end{tabularx}
\end{caseglance}

\begin{chapterquestions}
By the end of this case you will be able to answer these questions:
\begin{enumerate}[leftmargin=*,itemsep=2pt]
  \item Where should a language model sit in a fraud system, and where must it never sit?
  \item How do you offer a statistical anomaly model to an agent as an MCP\index{Model Context Protocol (MCP)} tool, with its score and its reasons?
  \item How do you set block and review thresholds when errors cost very different amounts and analysts are scarce?
  \item How do you check that an agent's explanation matches what the model actually used?
  \item How do you evaluate a fraud system when labels arrive months late and accuracy means nothing?
\end{enumerate}
\end{chapterquestions}

From this chapter on, cases are designed from a blank page. There is no sample to critique; the first alternative is the design an inexperienced team would most likely propose, and the method of Chapter~\ref{ch:toolkit} does the rest.

% ----------------------------------------------------------------------
\section{Business Brief: What Is Going Wrong Today?}
\label{sec:fr-brief}

Northgate issues credit cards to 3.2 million cardholders. About 600{,}000 card-not-present transactions---online and in-app purchases---arrive each day. When a merchant asks Northgate to authorize one, Northgate must answer within about 100 milliseconds. A fraud-scoring model, trained on years of labeled transactions, gives each one a score; rules then approve, decline, or approve and place the transaction in a queue for analysts to review after the fact.

Fourteen fraud analysts work each shift. For each queued transaction, an analyst opens the cardholder's history, compares the transaction with the cardholder's usual behavior, checks the merchant, decides, and, if fraud is likely, blocks the card and arranges a new one. A review takes about six minutes. The queue is always longer than the analysts' capacity; transactions not reached before the end of the shift are released unreviewed.

\begin{thinkabout}
Which mistake would you rather make: block a good customer's card, or let a fraudulent purchase through? How many good customers would you be willing to block to stop one fraud?
\end{thinkabout}

The economics are asymmetric in a way that dominates every design choice:
\begin{itemize}[itemsep=1pt]
  \item A \textbf{missed fraud} costs Northgate the transaction amount (for card-not-present fraud, the issuer usually bears it) plus about \$25 of chargeback handling.
  \item A \emph{wrongly declined or blocked good customer} costs much less per case---about \$6 in calls and goodwill---but there are many more good customers than fraudsters, and some of them stop using the card.
  \item An \textbf{analyst review} costs about \$0.70 a minute.
\end{itemize}

Fraud losses run at about 9 basis points (0.09\%) of card-not-present volume; the head of fraud wants them lower without declining more good customers, and the analysts' queue has grown 30\% in a year.

% ----------------------------------------------------------------------
\section{Stakeholders and KPIs: Who Cares, and How Will We Know?}
\label{sec:fr-kpis}

\begin{table}[htbp]
\centering
\small
\caption{Stakeholders of fraud review.}
\label{tab:fr-stakeholders}
\begin{tabularx}{\textwidth}{@{}>{\raggedright\arraybackslash}p{2.0cm}>{\raggedright\arraybackslash}X@{}}
\toprule
\rowcolor{coHead} Role & Who it is \\
\midrule
\cotag{coBlue}{Sponsor} & The head of fraud strategy. \\
\cotag{coGreen}{Users} & The fraud analysts. \\
\cotag{coAmber}{Approvers} & Model risk management, which validates every model used in decisions; compliance, for customer-contact rules; and information security. \\
\cotag{coPurple}{Owners} & The fraud analytics team (the scoring model), the authorization platform team, and fraud operations (the queue and the analysts). \\
\bottomrule
\end{tabularx}
\end{table}

\paragraph{KPIs and risks.} Table~\ref{tab:fr-kpis} lists them.
\begin{table}[htbp]
\centering
\small
\caption{KPIs (blue) and risks (red) for fraud review.}
\label{tab:fr-kpis}
\begin{tabularx}{\textwidth}{@{}>{\raggedright\arraybackslash}p{4.2cm}>{\raggedright\arraybackslash}X@{}}
\toprule
\rowcolor{coHead} Name & Definition and target \\
\midrule
\kpi{Fraud losses} & Fraud losses in basis points of card-not-present volume: down 20\%. \\
\kpi{Wrongly blocked} & Good customers wrongly declined or blocked per 10{,}000 transactions: no higher than today. \\
\kpi{Reviews per hour} & Transactions reviewed per analyst hour: up 50\%. \\
\risk{Model in hot path} & Any language model in the authorization path, where latency and availability are critical. \\
\risk{Unfaithful reason} & An explanation that misstates why a transaction was flagged, leading an analyst to a wrong decision. \\
\risk{Data exposure} & Cardholder personal data exposed beyond the analysts who need it. \\
\bottomrule
\end{tabularx}
\end{table}

\kpi{Fraud losses} and \kpi{Wrongly blocked} are a pair in the sense of Chapter~\ref{ch:business}: any design can cut fraud by declining more, and any design can please customers by declining less.

\begin{keyidea}[The value hypothesis, in one sentence]
If analysts receive each flagged transaction with a faithful summary of why the model flagged it and the context they would otherwise gather by hand, they will review substantially more transactions per hour; with more review capacity, fewer transactions need to be blocked outright, so fraud losses and wrongly blocked customers can fall together.
\end{keyidea}

% ----------------------------------------------------------------------
\section{Alternative Designs: What Could We Build?}
\label{sec:fr-designs}

\subsection{The agent-judge design: an agent that judges each transaction}

The design a team new to agents often proposes is to give a model each flagged transaction, the cardholder's history and some tools, and ask it whether the transaction is fraud. It fails step~3 of the design method\index{design method} on two of its four questions, judgment or rule and what must be known: the judgment ``is this fraud?'' is exactly what a statistical model trained on millions of labeled transactions does well, and what a language model reading a few dozen does poorly and expensively. If the agent-judge design is placed in the authorization path it also violates \risk{Model in hot path}: authorization needs an answer in 100 milliseconds, every time, which no agent loop\index{agent loop} can promise. We score the agent-judge design for completeness.

\subsection{The assistant design: the model scores, the agent assists the analyst}

The assistant design assigns each block to what it does best. Table~\ref{tab:fr-d2} writes it down as a design table\index{design table}.
\begin{itemize}[itemsep=1pt]
  \item The \textbf{scoring model} stays in the authorization path, served with low latency by the authorization platform. It is also exposed, read-only, through an MCP server\index{Model Context Protocol (MCP)!server} owned by fraud analytics, which returns for any transaction its score, the model version and the features that contributed most to the score (its \emph{attributions}).
  \item \textbf{Code} applies two thresholds: above the block threshold, the card is blocked at once; between the review and block thresholds, the transaction is approved and queued for review. Both thresholds are configuration, owned by fraud strategy and set by the analysis in Section~\ref{sec:fr-tradeoffs}.
  \item For each queued transaction, a \textbf{case agent} (Level~4, at most eight tool calls\index{tool call}) gathers what an analyst would: the attributions, the cardholder's recent transactions and devices, the merchant's fraud history, and similar confirmed cases. It writes a \textbf{case summary}: the facts, the model's top factors in plain language, and what is unusual against the cardholder's own history.
  \item A \textbf{faithfulness guardrail\index{guardrail}} checks that every reason the summary cites is one of the model's top factors (Section~\ref{sec:fr-build}). Summaries that fail are regenerated once, then shown with a warning.
  \item The \emph{analyst} decides, through the existing case tool. The agent has no write tools.
\end{itemize}

\begin{table}[htbp]
\centering
\small
\caption{The assistant design as a design table. Step 1 is the latency-critical path, and it contains no language model. The agent works on the review queue, after authorization, and has no write tools.}
\label{tab:fr-d2}
\begin{tabularx}{\textwidth}{@{}c>{\raggedright\arraybackslash}X>{\raggedright\arraybackslash}p{2.6cm}>{\raggedright\arraybackslash}p{3.4cm}>{\raggedright\arraybackslash}p{3.0cm}@{}}
\toprule
\rowcolor{coHead} \# & Step & Done by & Uses & Why it is there \\
\midrule
1 & Within 100 ms, score the transaction and apply the thresholds. Above the block threshold, block; in the review band, approve and queue; below it, approve & \tagCode\ No language model & Scoring model on the authorization platform; thresholds (owner: Fraud strategy) & \kpi{Fraud losses} \kpi{Wrongly blocked} \risk{Model in hot path} \\
2 & Order the review queue by expected loss & \tagCode & Review queue & \kpi{Fraud losses} \\
3 & Gather the context an analyst would and write a case summary & \tagAgent\ Level 4; at most eight calls; stops at the summary; no write tools & Scoring server (Fraud analytics): \texttt{explain\_score}, \texttt{similar\_cases}; card-history server (Cards team, read-only): \texttt{recent\_txns}, \texttt{devices}, \texttt{merchant\_profile} & \kpi{Reviews per hour} \risk{Data exposure} \\
4 & Check that every cited reason is one of the model's top factors. Regenerate once; then show with a warning & \tagGuardrail & Attributions from \texttt{explain\_score} & \risk{Unfaithful reason} \\
5 & Decide, through the existing case tool & \tagPerson\ Fraud analyst & The summary and the model's raw factors & \kpi{Fraud losses} \kpi{Wrongly blocked} \\
\bottomrule
\end{tabularx}
\end{table}

\subsection{The verification design: ask the cardholder}

The verification design adds a step to the assistant design: for transactions in a middle band, the system asks the cardholder through the bank's own app, with a push notification, ``Did you make this purchase?'' A ``yes'' clears the transaction without an analyst; a ``no'' blocks the card at once. The verification message is a fixed template sent through the authenticated app channel---never free text written by a model, and never by SMS or e-mail, which fraudsters imitate. The verification design needs compliance approval for customer contact and engineering work by the mobile team. Table~\ref{tab:fr-designs} compares the three designs.

\begin{table}[htbp]
\centering
\small
\caption{The three fraud designs side by side.}
\label{tab:fr-designs}
\begin{tabularx}{\textwidth}{@{}>{\raggedright\arraybackslash}p{2.7cm}>{\raggedright\arraybackslash}X>{\raggedright\arraybackslash}X>{\raggedright\arraybackslash}X@{}}
\toprule
\rowcolor{coHead} & \cellcolor{coRedBg}Agent judge & Assistant & Verification \\
\midrule
Who judges fraud & A language model & The statistical model, then an analyst & The same, plus the cardholder in a middle band \\
The 100-ms path & May contain the agent, which fails & Scoring model only & Scoring model only \\
The analyst & Checks the agent's verdict & Decides, helped by a summary & Decides what the cardholder does not resolve \\
Customer contact & None & None & A fixed push message in the bank's own app \\
\bottomrule
\end{tabularx}
\end{table}

% ----------------------------------------------------------------------
\section{Trade-offs: Which Design Should We Pilot?}
\label{sec:fr-tradeoffs}

\subsection{Thresholds under asymmetric costs and limited capacity}

The business skill of this case is choosing thresholds when errors cost different amounts and the review budget is fixed. Listing~\ref{lst:fraud-policy} simulates a day of card-not-present transactions with a realistic, imperfect score, and searches for the block and review thresholds that minimize total daily cost: missed fraud, wrongly blocked customers and analyst time. Analysts work the review queue from the highest score down until their capacity runs out.

\lstinputlisting[language=Python, caption={Block and review thresholds under capacity (\texttt{ch08/fraud\_policy.py}, core).}, label={lst:fraud-policy}, firstline=6, lastline=36]{code/fraud_policy.py}

\begin{thinkabout}
Before reading the output, predict: if summaries let analysts review faster, will the main benefit be more fraud caught, or fewer good customers blocked?
\end{thinkabout}

\begin{lstlisting}[language={}, numbers=none, frame=none, xleftmargin=0em, basicstyle=\ttfamily\scriptsize]
fraud cases 734, fraud value $100,860
analyst alone (6 min/alert)        review>=0.05 block>=0.55  reviewed  1050  caught 476/734  false blocks  839  cost $52,994/day
with case summary (3.5 min/alert)  review>=0.05 block>=0.60  reviewed  1800  caught 482/734  false blocks  294  cost $48,533/day
no review: block >= 0.70           cost $85,681/day
\end{lstlisting}

Four things in this output generalize.
\begin{itemize}[itemsep=2pt]
  \item \emph{Capacity, not the review threshold, decides what is reviewed.} Because the queue is worked from the top and always exceeds capacity, the optimal review threshold sinks to the bottom of the search range; what matters is how many items analysts can reach. That is why the case agent targets review time.
  \item \emph{Faster review moves the block threshold.} With case summaries, analysts review 1{,}800 transactions instead of 1{,}050, so fewer need to be blocked blindly: the optimal block threshold rises from 0.55 to 0.60 and wrongly blocked customers fall from 839 to 294 a day. \kpi{Wrongly blocked} improves because of an assistant to \kpi{Reviews per hour}---a connection that only the joint optimization reveals.
  \item \emph{The asymmetry sets the shape.} Because a missed fraud costs about \$160 on average and a wrongly blocked customer about \$6, the optimizer blocks aggressively when review capacity is short. Change the friction cost and the optimum moves; Exercise~4 asks you to find where.
  \item \emph{The model's quality is the ceiling.} Even the best policy catches only about two-thirds of fraud cases, because the simulated score is deliberately imperfect. No amount of agent assistance changes that; a better scoring model does.
\end{itemize}

\begin{keyidea}
The agent is not a substitute for the model. It makes people faster; only a better model makes the fraud signal better.
\end{keyidea}

At Northgate's volumes, the difference between the first two rows is about \$4{,}500 a day, or \$1.6 million a year, before counting the retention value of 545 fewer wrongly blocked customers each day.

\subsection{Scores and decision}

\begin{table}[htbp]
\centering
\small
\caption{Trade-off table for fraud review. The agent-judge design is shown outside the authorization path; inside it, the agent-judge design fails \risk{Model in hot path}.}
\label{tab:fr-tradeoff}
\begin{tabular}{@{}lcccc@{}}
\toprule
\rowcolor{coHead} & Weight & Agent judge & Assistant & Verification \\
\midrule
Quality (fraud caught) & 0.30 & 2 & 4 & 5 \\
Customer impact & 0.20 & 2 & 4 & 4 \\
Speed and latency & 0.10 & 2 & 5 & 4 \\
Security and privacy & 0.15 & 2 & 4 & 3 \\
Control and auditability & 0.15 & 1 & 5 & 4 \\
Maintainability & 0.10 & 3 & 4 & 3 \\
\midrule
\rowcolor{coAmberBg} Weighted score & 1.00 & 1.95 & 4.25 & 4.05 \\
\bottomrule
\end{tabular}
\end{table}

The verification design catches more fraud, because many mid-band transactions are resolved by the cardholder in seconds. It scores lower on privacy and control because it introduces a new customer-facing channel with its own failure modes, and it depends on another team and on compliance approval. The decision is the assistant design now, with the verification design as a separate project once the assistant design's evaluation establishes the baseline it would improve on.
